
%% USPSC-Cap3-Conclusao.tex
% ---
% Conclusão
% ---
\chapter{Conclusion}\label{cap:conclusion}

The main question investigated in this work was how quantum properties affect the heating-cooling asymmetry in discrete and continuous-variable systems. The relaxation dynamics was described by the Lindblad master equations and we used the thermal kinematics approach to identify the heating-cooling asymmetry. In this respect, we introduced an adaptation of thermal kinematics to Gaussian systems. In addition to thermal kinematics, we analysed how entropy production can be used to explain which thermal relaxation process is faster.

For the discrete case, the two-interacting-qubits model allowed us to study the thermal relaxation of systems with coherence and entanglement. By applying Davies maps to this model, the global state always thermalizes to the bath's temperature, but the interaction between the qubits can generate correlations between them, and, as a result, we cannot write the global steady state as the tensor product of qubit's thermal states. However, we can associate an effective temperature to the qubits with the cost of this temperature at equilibrium being lower than the bath's temperature. Another consequence of the qubits interaction generating correlations is that the local entropy production rate becomes negative at some periods, which indicates the a local non-Markovian dynamics. We noticed that the thermal relaxation case with non-Markovianity was faster than the no-interaction and full Markovian case. However, we observed that strong interaction causes entanglement, which reduces the asymmetry and make the relaxation slower than the weaker interaction without entanglement generation. Additionally, we saw that initial global coherence can apparently increase the asymmetry, changing the intermediary dynamics, but it do not affect the total equilibration time of the qubits.

For the continuous-variable case, we explored the heating-cooling asymmetry in single-mode Gaussian states. We considered that squeezing can generate quantum resources, such as coherences, in Gaussian states. However, instead of directly quantifying these resources, we used ergotropy to separate how much of the nonequilibrium energy is thermodynamically useful. Because the passive state of Gaussian systems is always a thermal state, any quantumness of the system will be related to ergotropic quantities. An interesting consequence of using the ergotropic approach is the possibility of separating the Wigner-Fisher information and the entropy production rate into passive and ergotropic contributions. So, we can separate the asymmetry origin from passive contributions (that have no quantumness for Gaussian systems) and from ergotropic contribution (that may have quantumness).

The first Gaussian system we analysed was thermal states, in which we analytically showed that heating is always faster than cooling and where entropy production successfully explained its origin. Second, we compared the relaxation of different thermal states that were equally displaced and showed that displacement does not affect the asymmetry, being the passive contribution its generator. Again, the total entropy production explained heating being faster than cooling and separating the contributions showed that the ergotropy was the same for both heating and cooling, which is a consequence of being independent of the system's initial temperature. Then, we study the asymmetry in squeezed thermal states where both ergotropic and passive contributions are affected by the Gaussian operation. In this case, using bigger squeezing parameters resulted in a reduction of the asymmetry of the total state while the passive state had its asymmetry increased. To explain this apparent contradiction, we observed that higher squeezing parameter reduced the initial entropy production of the hot and cold passive states, but it increase their difference. However, we noticed that increasing the initial ergotropy did not lead to enough increase in the ergotropic contribution of entropy production to compensate the passive reduction. As a result, the difference of the heating and cooling total entropy production decreases, causing the asymmetry reduction. A last result from the asymmetry with squeezed states it that we identified a non-Markovian dynamics for the heating passive state of the big squeezing state and we also observed an acceleration of the dynamics. As a result, once again we have an indication that non-Markovian dynamics can help a faster relaxation. Until this moment, all results showed that heating was faster than cooling, but it is known from literature that this asymmetry can be inverted. This motivated the investigation of how we could have cooling faster than heating in our systems. Using the knowledge that squeezed states relaxate faster than displaced states, we squeezed the hotter state and displaced the colder state. With this combination we inverted the asymmetry partially (for smaller parameters) and totally (for bigger parameters). We learned that the full inversion is generated only by the ergotropic contribution, while heating is still faster than cooling for the passive states. Also, we discussed that the partial inversion is actually a consequence of squeezing having a faster relaxation than displaced states, which was already discussed in the ergotropic Mpemba effect \cite{medina-2025}. In this case, the initial entropy production failed to indicate the faster process, but its time evolution helped understand the partial inversion. 

Beyond the investigations done in this work, there is a variety of quantum properties and systems where heating-cooling asymmetry could be generated. Regarding the same model of two-interacting qubits, there is also the possibility of investigating other interaction hamiltonians, such as XXX, XXZ, XYZ, and transverse field Ising. In the context of Gaussian states, the ergotropic approach was not a clear indicator of quantumness. So, a next step can be quantifying the quantum resources, such as coherence, to understand how much of the behaviour we studied is actually a consequence of quantum properties. In this respect, extending the analysis for more modes can allow the investigation of the heating-cooling asymmetry in entangled Gaussian systems. As we briefly discussed in Sec. \ref{sec:result-squeezed}, the passive state of squeezed states can have a non-Markovian dynamics that seems to accelerate the relaxation. So, further investigating this aspect can also be an interesting direction. Another possible future direction would be to study why entropy production is asymmetric. A possible direction is using the Onsager relations for far-from-equilibrium Gaussian states as derived in \cite{landi-2020}. The Onsager relations write the entropy production in terms of currents and Ref. \cite{landi-2020} showed that it becomes non-linear and a not even function for arbitrarily far-from-equilibrium Gaussian states.

